%% LyX 2.4.4 created this file.  For more info, see https://www.lyx.org/.
%% Do not edit unless you really know what you are doing.
\documentclass[english]{article}
\usepackage[T1]{fontenc}
\usepackage[latin9]{inputenc}
\usepackage{enumitem}
\usepackage{amsmath}
\usepackage{amsthm}
\usepackage{geometry}
\geometry{verbose,tmargin=1in,bmargin=1in,lmargin=1in,rmargin=1in}

\makeatletter
%%%%%%%%%%%%%%%%%%%%%%%%%%%%%% Textclass specific LaTeX commands.
\numberwithin{equation}{section}
\numberwithin{figure}{section}
\newlength{\lyxlabelwidth}      % auxiliary length 

%%%%%%%%%%%%%%%%%%%%%%%%%%%%%% User specified LaTeX commands.
\usepackage{fancyhdr}
\usepackage{lastpage}
\pagestyle{fancy}
\usepackage{enumitem}
\usepackage{ifthen}
\everymath=\expandafter{\the\everymath\displaystyle}
%\DeclareMathSizes{10}{10}{10}{10}
\usepackage{multicol}

\usepackage{tikz}
\usepackage{tikz-3dplot}
\usetikzlibrary{calc,arrows}

\usetikzlibrary{patterns}
\usetikzlibrary{plotmarks}
\usepackage{pgfplots}
\pgfplotsset{compat=newest}

\usepackage{calc}
\usepackage[nomessages]{fp}

\usepackage{datetime}
% Added by lyx2lyx
\setlength{\parskip}{\smallskipamount}
\setlength{\parindent}{0pt}

\makeatother

\usepackage{babel}
\begin{document}
\renewcommand{\author}{Dr.\,James Ashe}

\newcommand{\classNum}{336}

\newcommand{\classSec}{A}

\newcommand{\examType}{HW 1.3 Solutions}

\renewcommand{\date}{\formatdate{6}{9}{2013}}

%%First page's header%%%%%%%%%%%%%%%%%%%%%%
\fancypagestyle{firststyle} %%header settings for first pagr
{
	\fancyhead[L]{MTH \classNum{}(\classSec{}) \author{}\\
	\textbf{\examType{}} ~\date{}}
	\fancyhead[C]{}
	\fancyhead[R]{\textbf{Name:\hspace{4cm}}}
	\setlength{\headsep}{14pt}
}
%%%%%%%%%%%%%%%%%%%%%%%%%%%%%%%%%%%%%%%%%%

%%Next page's header%%%%%%%%%%%%%%%%%%%%%%%
\setlist{nolistsep}
\lhead{\classNum{}(\classSec{})} 
\chead{\textbf{Name:}} 
\cfoot{\thepage{}~of~\pageref{LastPage}} 
\setlength{\headsep}{5pt} 
%%%%%%%%%%%%%%%%%%%%%%%%%%%%%%%%%%%%%%%%%%%

%%Various settings; see the preamble for other settings%%%%%
%\setenumerate[0]{leftmargin=12pt,itemindent=0pt,labelwidth=0pt}
\setlist{topsep=.5pt}
\renewcommand{\labelenumi}{\textbf{\arabic{enumi}.}}
\renewcommand{\labelenumii}{\textbf{\alph{enumii}.}}
\thispagestyle{firststyle}
%%%%%%%%%%%%%%%%%%%%%%%%%%%%%%%%%%%%%%%%%%

\newcommand{\xmax}{0}
\newcommand{\xmin}{-\xmax}
\newcommand{\xstep}{0}
\newcommand{\ymax}{0}
\newcommand{\ymin}{-\ymax}
\newcommand{\ystep}{0} 
\newcommand{\dspace}{\vspace{1cm}}
\newcommand{\xa}{0}
\newcommand{\xb}{0}
\newcommand{\ya}{1}
\newcommand{\yb}{1}
\begin{enumerate}[resume]
\item In general, given a normal vector $\mathbf{a}$ for a hyperplane,
$\mathbf{a\cdot}(\mathbf{x}-\mathbf{x}_{0})=0$ for any two points
$\mathbf{x}$ and \textbf{$\mathbf{x}_{0}$ }on the plane (since $\mathbf{x}-\mathbf{x}_{0}$
is a vector on the plane and hence perpendicular to $\mathbf{a}$).
So $\mathbf{a}\cdot\mathbf{x}=\mathbf{a}\cdot\mathbf{x}_{0}$. If
we let $\mathbf{x=}(x_{1},x_{2},x_{3})$ then $\mathbf{a}\cdot\mathbf{x}=\mathbf{a}\cdot\mathbf{x}_{0}$
gives a general solution when we know a single point on the plane
$\mathbf{x}_{0}$.

\begin{enumerate}[resume]
\item The direction vector is $(3,2)$. So a normal vector $(-2,3)$ and
the point $(-1,2)$ gives the equation 
\begin{eqnarray*}
(-2,3)\cdot\mathbf{x} & = & (-2,3)\cdot(-1,2)=8\\
\Rightarrow(-2,3)\cdot(x_{1},x_{2},x_{3}) & = & \underline{-2x_{1}+3x_{2}=8}
\end{eqnarray*}
\item [\textbf{c.}]\setcounter{enumii}{3}A normal vector to the plane is
$(1,2,-2)$ (given by the direction vector of the orthogonal line).
Using the given point $(2,0,1)$ we have,
\begin{eqnarray*}
(1,2,-2)\cdot\mathbf{x} & = & (1,2,-2)\cdot(2,0,1)=0\\
\Rightarrow(1,2,-2)\cdot(x_{1},x_{2},x_{3}) & = & \underline{x_{1}+2x_{2}-2x_{3}=0}
\end{eqnarray*}
\item A normal vector will be will be orthogonal to both $(1,1,1)$ and
$(2,1,0)$. So then $\mathbf{a}\cdot(1,1,1)=0\Rightarrow a_{1}+a_{2}+a_{3}=0$
and $\mathbf{a\cdot}(2,1,0)=0\Rightarrow2a_{1}+a_{2}=0$. Substituting
$a_{2}=-2a_{1}$ into the first equation we get $a_{3}=a_{1}$. So
then $\mathbf{a=}(1,-2,1)$ works. A Cartesian equation is then 
\[
(1,-2,1)\cdot\mathbf{x}=(1,-2,1)\cdot(1,1,2)=1\Rightarrow\underline{x_{1}-2x_{2}+x_{3}=1}
\]
\end{enumerate}
\item [\textbf{3.}]\vspace{0cm}

\begin{enumerate}
\item $x_{1}-2x_{2}+3x_{3}=4\Rightarrow x_{1}=2x_{2}-3x_{3}+4$. So then
\begin{eqnarray*}
\mathbf{x} & = & (2x_{2}-3x_{3}+4,x_{2},x_{3})\mbox{.}\\
 & = & \underline{(4,0,0)+x_{2}(2,1,0)+x_{3}(-3,0,1)}
\end{eqnarray*}
\item [\textbf{c.}]$x_{1}+x_{2}-x_{3}+2x_{4}=5\Rightarrow x_{1}=5-x_{2}+x_{3}-2x_{4}$.
So then 
\begin{eqnarray*}
\mathbf{x} & = & (5-x_{2}+x_{3}-2x_{4},x_{2},x_{3},x_{4})\\
 & = & \underline{(5,0,0,0)+x_{2}(-1,1,0,0)+x_{3}(1,0,1,0)+x_{4}(-2,0,0,1)}
\end{eqnarray*}
\end{enumerate}
\item [\textbf{5.}]Alternatively to the methods shown, you could use the
techniques discussed in 1.4.

\begin{enumerate}
\item We have that $x_{1}+x_{2}+x_{3}=1\Rightarrow x_{1}=1-x_{2}-x_{3}$
substituting this into the other equation we have, $2(1-x_{2}-x_{3})+x_{2}+2x_{3}=1\Rightarrow x_{2}=1$.
By substituting this into the first equation we see that $x_{2}=1\Rightarrow x_{1}=-x_{3}$.
So then we have,
\begin{eqnarray*}
\mathbf{x} & = & (x_{1},1,-x_{1})\\
 & = & \underline{(0,-1,0)+x_{1}(1,0,-1)}
\end{eqnarray*}
\item $x_{1}-x_{2}=1\Rightarrow x_{1}=1+x_{2}$ substituting this into the
other equation we have. $(1+x_{2})+x_{2}+2x_{3}=5\Rightarrow x_{2}=2-x_{3}$.
By substituting this into the first equation we see that $x_{1}=3-x_{3}$.
Now all the variables are expressed in terms of $x_{3}$ so we have,
\begin{eqnarray*}
\mathbf{x} & = & (3-x_{3},2-x_{3},x_{3})\\
 & = & \underline{(3,2,0)+x_{3}(-1,-1,1)}
\end{eqnarray*}
Alternatively you might have solved for $x_{2}$ and gotten $\mathbf{x}=(1,0,2)+x_{2}(1,1,-1)$.
It is also typical to write the scalar as $t$ or $s$ instead of
$x_{i}$ so that it would not be confused as a dependent parameter. 
\end{enumerate}
\end{enumerate}

\end{document}
